% !TEX TS-program = lualatexmk
% !TEX parameter =  --shell-escape

\documentclass[vespers_book_la_eng.tex]{subfiles}

\ifcsname preamble@file\endcsname
  \setcounter{page}{\getpagerefnumber{M-vespers_book_la_eng_sundays_advent}}
\fi

\begin{document}

\phantomsection
\chapter*{SUNDAYS OF ADVENT.}\header{sundays of advent.}
\addcontentsline{toc}{chapter}{Sundays of Advent.}

% \twosided[pc] %% to allow for the left-right switching of columns needed for bilingual liturgical typesetting (like in missals); it's unclear why this goes BEFORE everything else and not when activating the environment. %% It's OK for paper booklets to not switch, but it's probably best practice to switch.
% 

\bigtitle{I Sunday of Advent.}
\addcontentsline{toc}{section}{I Sunday of Advent} \phantomsection

\firstantiphon{8. G}

\gscore[]{}{In_illa_die_incipit}

\psalmus{109}{109_8G}{109_8}{109_en}

\smallscore{an_in_illa_die_solesmes_1961}

\antiphona{english}{In that day the mountains shall drop down sweet wine, and the hills shall flow with milk and honey. Alleluia}

\gscore[2. Ant.]{8. G}{Juncundare_incipit}

\psalmus{110}{110_8Gstar}{110_8}{110_en}

\smallscore{an_jucundare_filia_solesmes_1961}

\antiphona{english}{Sing, O daughter of Zion, and rejoice with all the heart, O daughter of Jerusalem. Alleluia}

\gscore[3. Ant.]{5. a}{Ecce_Dominus_veniet_incipit}

\psalmus{111}{111_5a}{111_5}{111_en}

\smallscore{an_ecce_dominus_veniet_solesmes_1961}

\antiphona{english}{Behold, the Lord shall come, and all his saints with him; and it shall come to pass in that day that the light shall be great. Alleluia}

\gscore[4. Ant.]{7. c}{Omnes_sitientes_incipit}

\psalmus{112}{112_7c}{112_7}{112_en}

\smallscore{an_omnes_sitientes_solesmes_1961}

\antiphona{english}{Lo, every one that thirsteth come ye to the waters: seek ye the Lord while He may be found. Alleluia}

\gscore[5. Ant.]{4. A*}{Ecce_veniet_incipit}

\psalmus{113}{113_4_alt_Astar}{113_4A_A_star}{113_en}

\smallscore{an_ecce_veniet_propheta_solesmes_1961}

\antiphona{english}{Behold, a great Prophet shall arise, and He shall build up a new Jerusalem. Alleluia}

\rubrique{\begin{enpars}
The chapter is proper to each Sunday of Advent, while the hymn and verse are always the same on the Sundays of the same season.\end{enpars}}

\chapterlateng
\scripture{Rom. 13, 11.}

\texteparacol{\lettrine{F}{r}atres: Hora est jam nos de somno súrgere:~† nunc enim própior est nostra salus,~* quam cum credídimus.

\rr Deo grátias.}{english}{\noindent And that knowing the season; that it is now the hour for us to rise from sleep. For now our salvation is nearer than when we believed. Thanks be to God.}

\rubrique{The response made at the end of the chapter is always the same.}
%% The usual tone and the 3rd tone, used ad libitum, are omitted as we do not use these tones, but they are included in the tex file and the gabc files are available, for the sake of completeness.

\hymnlateng
\label{creator_alme} 
\gscore[]{4.}{hy_creator_alme_siderum_solesmes_1961}

\textes{creator_alme_siderum_en}
%\rubrique{\begin{enpars}
%The cantor sings the verse to which all reply.
%\end{enpars}}
\label{versiculus_adventus} 
\smallscore{versiculus_adventus}

%\gresetinitiallines{0}
%
%\gabcsnippet{(c3)℣. Ro(h)rá(h)te(h) cæ(h)li(h) dé(h)su(h)per,(h) et(h) nu(h)bes(h) plu(h)ant(h) ju(h)stum.(g_hvGFEfgf.) (::)}
%
%\gabcsnippet{(c3) ℟. A(h)pe(h)ri(h)á(h)tur(h) ter(h)ra,(h) et(h) gér(h)mi(h)net(h) Sal(h)va(h)tó(h)rem.(g_hvGFEfgf.) (::)}

\begin{enpars} \noindent \vv Drop down dew, ye heavens, from above, and let the clouds rain the Just.

\noindent \rr Let the earth be opened, and bud forth the Saviour.
 \end{enpars}

%\gscore[Ad Magnif.]{Ant. 8. G}{Ne_timeas_Maria_incipit}

\gscore[Ad Magnif.]{Ant. 8. G}{an_ne_timeas_solesmes_1961}

\antiphona{english}{Fear not, Mary, for thou hast found grace with the Lord; behold, thou shalt conceive in thy womb and bring forth a son. Alleluia}

\collectlateng
\texteparacol{
\lettrine{E}{x}cita, quǽsumus Dómine, poténtiam tuam, et veni:~† ut ab imminéntibus peccatórum nostrórum perículis, te mereámur protegénte éripi,~* te liberánte salvári. Qui vivis et regnas cum Deo Patre.}{english}{\noindent Let us pray. Stir up, O Lord, we pray thee, thy strength, and come among us, that whereas through our sins and wickedness we do justly apprehend thy wrathful judgments hanging over us, thy bountiful grace and mercy may speedily help and deliver us; Who livest and reignest with God the Father.}

 \rubrique{\begin{enpars}The Suffrage of All Saints is omitted for all of Advent.\end{enpars}}

%\rubrique{\begin{enpars}
%The cantor sings the verse to which all reply.
%\end{enpars}}

%\gscore[]{6.}{BD_Dominicis_Adventus_Quadragesimae}

%\vv Let us bless the Lord.
%
%\rr Thanks be to God.
%
%\textebilingue{fidelium_animae}{fidelium_animae_en}

%\rubrique{\begin{enpars}
%Then one Our Father is said, unless Compline follows, or even some pious exercises, a sermon, or Benediction of the Blessed Sacrament, in parish and other smaller churches.\end{enpars}}

%\rubrique{\begin{enpars}
%Then the office is concluded with the following verses, the Marian antiphon, etc.
%\end{enpars}}
\bigtitle{II Sunday of Advent.}
\addcontentsline{toc}{section}{II Sunday of Advent} \phantomsection

\firstantiphon{1. g}

\gabcsnippet{(c4) EC(dd)ce(c) in(d) nú(e_[uh:l]f)bi(g)bus(fe) cæ(d)li(d.c.) (::)}

\psalmus{109}{109_1g}{109_1}{109_en}

\smallscore{an_ecce_in_nubibus_solesmes_1961}

\antiphona{english}{Behold, the Lord cometh in the clouds of heaven with great power. Alleluia}

\greannotation{2. Ant.}
\greannotation{7. d}
\gabcsnippet{(c2) U{R}bs(cg..) (::)}

\psalmus{110}{110_7d}{110_7}{110_en}

\smallscore{an_urbs_fortitudinis_solesmes_1961}

\antiphona{english}{Our Zion is a strong city, the Saviour will God appoint in her for walls and bulwarks; open ye the gates, for God is with us. Alleluia}

\greannotation{3. Ant.}
\greannotation{7. a}
\gabcsnippet{(c3) EC(e')ce(e) ap(g')pa(h)ré(ih/ij)bit(i.) (::)}

\psalmus{111}{111_7a}{111_7}{111_en}

\smallscore{an_ecce_apparebit_solesmes_1961}

\antiphona{english}{Behold, the Lord shall appear and not lie though He tarry, wait for him, because He will come and will not tarry. Alleluia}

\greannotation{4. Ant.}
\greannotation{1. f}
\gabcsnippet{(c4) MOn(d_c~)tes(f) et(g) col(f_h)les(h.) (::)}

\psalmus{112}{112_1f}{112_1}{112_en}

\smallscore{an_montes_et_colles_solesmes_1961}

\antiphona{english}{The mountains and the hills shall break forth before God into singing, and all the trees of the wood shall clap their hands for the Lord the Ruler cometh, and He shall reign for ever and ever. Alleluia, Alleluia}

\greannotation{5. Ant.}
\greannotation{3. a}
\gabcsnippet{(c4) EC(g)ce(g) Dó(hi)mi(i)nus(i) no(i)ster(i.) (::)}

\psalmus{113}{113_3a}{113_3a_b}{113_en}

\smallscore{an_ecce_dominus_noster_ut_illuminet_solesmes_1961}

\antiphona{english}{Behold, our Lord cometh with power, and He shall lighten the eyes of his servants. Alleluia}

\chapterlateng
\scripture{Rom. 15, 4.}

\texteparacol{\lettrine{F}{r}atres: Quæcúmque scripta sunt, ad nostram doctrínam scripta sunt:~† ut per patiéntiam, et consolatiónem Scripturárum~* spem habeámus.}{english}{\noindent Brethren: For what things soever were written, were written for our learning: that through patience and the comfort of the scriptures, we might have hope.}

\rubrique{Hymnus \normaltext{Creátor alme síderum, \pageref{creator_alme}.} \normaltext{℣. Mane nobíscum, \pageref{versiculus_adventus}.}}

\gscore[Ad Magnif.]{8. G*}{an_tu_es_qui_venturus_es_an_alium_solesmes_1961}
\antiphona{english}{Art thou he that art to come or look we for another? Go and relate to John what you have seen: the blind see, the dead rise again, to the poor the gospel is preached. Alleluia}

\collectlateng

\texteparacol{
\lettrine{E}{x}cita Dómine corda nostra ad præparándas Unigéniti tui vias:~† ut per ejus advéntum,~* purificátis tibi méntibus servíre mereámur.
Qui tecum vivit et regnat.}{english}{\noindent Let us pray. Stir up our hearts, O Lord, to make ready the ways of thine Only-begotten Son, that by his coming our minds being purified, we may the more worthily give up ourselves to thy service: Who with thee liveth and reigneth.}

\newpage
\bigtitle{III Sunday of Advent.}
\addcontentsline{toc}{section}{III Sunday of Advent} \phantomsection

\firstantiphon{1. a}

\gscore[]{}{veniet_dominus_incipit}

\psalmus{109}{109_1a}{109_1}{109_en}

\smallscore{an_veniet_dominus_et_non_solesmes_1961}

\antiphona{english}{The Lord will come, and will not tarry; He both will bring to light the hidden things of darkness, and will make himself manifest to all people. Alleluia}

\gscore[2. Ant.]{7. b}{jerusalem_gaude_incipit}

\psalmus{110}{110_7b}{110_7}{110_en}

\smallscore{an_jerusalem_gaude_solesmes_1961}

\antiphona{english}{Rejoice greatly, O Jerusalem, for thy Saviour cometh unto thee. Alleluia}

\gscore[3. Ant.]{8. g}{dabo_incipit}

\psalmus{111}{111_8G}{111_8}{111_en}

\smallscore{an_dabo_in_sion_solesmes_1961}

\antiphona{english}{I will place salvation in Zion, and My glory in Jerusalem. Alleluia}

\gscore[4. Ant.]{5. a}{montes_et_omnes_incipit}

\psalmus{112}{112_5a}{112_5}{112_en}

\smallscore{an_montes_et_omnes_solesmes_1961}

\antiphona{english}{Every mountain and hill shall be made low, and the crooked shall be made straight, and the rough places plain; O Lord, come and make no tarrying. Alleluia}

\gscore[5. Ant.]{2. D}{juste_et_pie_incipit}

\psalmus{113}{113_2}{113_2}{113_en}

\smallscore{an_juste_et_pie_solesmes_1961}

\antiphona{english}{We should live righteously and godly, looking for that blessed hope and the coming of the Lord}

\chapterlateng

\scripture{Phil. 4, 4 – 5.}

\texteparacol{\lettrine{F}{r}atres: Gaudéte in Dómino semper: íterum dico, gaudéte.~† Modéstia vestra nota sit ómnibus homínibus:~* Dóminus enim prope est.}{english}{\noindent Brethren: Rejoice in the Lord always; again, I say, rejoice. Let your modesty be known to all men. The Lord is nigh.}

\rubrique{Hymnus \normaltext{Creátor alme síderum, \pageref{creator_alme}.} \normaltext{℣. Mane nobíscum, \pageref{versiculus_adventus}.}}

\gscore[Ad Magnif.]{8. G}{an_beata_es_maria_solesmes}

\antiphona{english}{Blessed art thou, O Mary, who hast believed the Lord, for there shall be a performance of those things which were told thee from the Lord. Alleluia}

\collectlateng

\texteparacol{\lettrine{A}{u}rem tuam quǽsumus Dómine, précibus nostris accómmoda:~† et mentis nostræ ténebras~* grátia tuæ visitatiónis illústra. Qui vivis et regnas cum Deo Patre.}{english}{\noindent Let us pray. O Lord, we beseech thee, mercifully incline thine ears unto our prayers, and lighten the darkness of our minds by the grace of thy heavenly visitation.
Who livest and reignest with God the Father.}

\section*{THE GREATER ANTIPHONS.}
\addcontentsline{toc}{section}{Greater Antiphons} \phantomsection

\label{o_antiphons} \phantomsection
 \rubrique{\begin{enpars}
The greater antiphon, sometimes called the \normaltext{O} antiphon, is sung in full before and after, or ``doubled'', standing both before and after the canticle, and from now until the end of the office.\end{enpars}}

 \rubrique{\begin{enpars}
\end{enpars}}

\bigtitle{December 17.}

\gscore[]{2. D}{an_o_sapientia_solesmes_1961}
\antiphona{english}{O Wisdom, that comest out of the mouth of the Most High, that reachest from one end to another, and dost mightily and sweetly order all things; come, to teach us the way of prudence}

\rubrique{Cant. \normaltext{Magníficat} in the solemn tone, \pageref{M-Mag_2_sol}.} %%pageref to that section since it's an exception to the rules

\rubrique{For the commemoration:}

\textes{versiculus_in_Adventu}

\bigtitle{December 18.}

\gscore[]{2. D}{an_o_adonai_solesmes_1961}
\antiphona{english}{O Adonai, and Ruler of the house of Israel, who didst appear unto Moses in the burning bush, and gavest him the law in Sinai; come, to redeem us with an outstretched arm}

\bigtitle{December 19.}

\gscore[]{2. D}{an_o_radix_jesse_solesmes_1961}
\antiphona{english}{O Root of Jesse, which standest for an ensign of the people, at whom the kings shall shut their mouths, to whom the Gentiles shall seek; come to deliver us, make no tarrying}
\label{an_o_clavis_david_solesmes_1961}

\bigtitle{December 20.}

\gscore[]{2. D}{an_o_clavis_david_solesmes_1961}
\antiphona{english}{O Key of David, and Sceptre of the house of Israel that openest, and no man shutteth; and shuttest and no man openeth; come to bring out the prisoners from the prison, and them that sit in darkness, and in the shadow of death}

\bigtitle{December 21.}
\label{an_o_oriens_solesmes_1961}
\gscore[]{2. D}{an_o_oriens_solesmes_1961}
\antiphona{english}{O Day-Spring, brightness of the everlasting Light, Sun of Righteousness; come, to give light to them that sit in darkness, and in the shadow of death}

\bigtitle{December 22.}

\gscore[]{2. D}{an_o_rex_gentium_solesmes_1961}
\antiphona{english}{O King of the Gentiles, yea, and Desire thereof, O Cornerstone that makest of twain one; come to save man, whom thou hast made of the dust of the earth}

\bigtitle{December 23.}

\gscore[]{2. D}{an_o_emmanuel_solesmes_1961}
\antiphona{english}{O Emmanuel, our King and our Law-giver, Longing of the Gentiles, yea, and Salvation thereof; come to save us, O Lord our God}

\bigtitle{IV Sunday of Advent.}

\addcontentsline{toc}{section}{IV Sunday of Advent} \phantomsection

\firstantiphon{1. g}

\gscore[]{}{canite_tuba_incipit}

\psalmus{109}{109_1g}{109_1}{109_en}

\smallscore{an_canite_tuba_solesmes_1961}

\antiphona{english}{Blow ye the trumpet in Zion, for the day of the Lord is nigh at hand: behold, He cometh to save us! Alleluia, Alleluia}

\gscore[2. Ant.]{1. f}{ecce_veniet_desideratus_incipit}

\psalmus{110}{110_1f}{110_1}{110_en}

\smallscore{an_ecce_veniet_desideratus_solesmes_1961}

\antiphona{english}{Behold, the desire of all nations shall come; and the house of the Lord shall be filled with glory. Alleluia}

\gscore[3. Ant.]{1. g}{erunt_prava_incipit}

\psalmus{111}{111_1g}{111_1}{111_en}

\smallscore{an_erunt_prava_solesmes_1961}

\antiphona{english}{The crooked shall be made straight, and the rough places plain; O Lord, come, and make no tarrying. Alleluia}

\gscore[4. Ant.]{1. f}{dominus_veniet_incipit}

\psalmus{112}{112_1f}{112_1}{112_en}

\smallscore{an_dominus_veniet_solesmes_1961}

\antiphona{english}{The Lord cometh! Go ye out to meet him, and say: How great is his dominion, and of his kingdom there shall be no end! He is the mighty God, the Ruler, the Prince of Peace. Alleluia, Alleluia}

\gscore[5. Ant.]{2. D}{omnipotens_sermo_incipit}

\psalmus{113}{113_2}{113_2}{113_en}

\smallscore{an_omnipotens_sermo_solesmes_1961}

\antiphona{english}{Thine Almighty Word, O Lord, shall leap down out of thy royal throne. Alleluia}

\chapterlateng
\scripture{1 Cor. 4, 1 – 2.}
%
\texteparacol{\lettrine{F}{r}atres: Sic nos exístimet homo ut minístros Christi,~† et dispensatóres mysteriórum Dei.~* Hic jam quǽritur inter dispensatóres, ut fidélis quis inveniátur.}{english}{\noindent Brethren: Let a man so account of us as of the ministers of Christ, and the dispensers of the mysteries of God. Here now it is required among the dispensers, that a man be found faithful.}

\rubrique{Hymnus \normaltext{Creátor alme síderum, \pageref{creator_alme}.} \normaltext{℣. Mane nobíscum, \pageref{versiculus_adventus}.}}

 \rubrique{\begin{otherlanguage}{english}
The \normaltext{O} antiphon corresponding to the day is sung, as above, \normaltext{\pageref{o_antiphons}.}, as there is no antiphon for the IV Sunday of Advent, which cannot fall earlier than December 18.\end{otherlanguage}}

%\rubrique{The ``O'' antiphon of the day is taken as the commemoration of Advent at Vespers of Saint Thomas, Apostle.}

\collectlateng

\texteparacol{
\lettrine{E}{x}cita, quǽsumus Dómine, poténtiam tuam, et veni: et magna nobis virtúte succúrre:~† ut per auxílium grátiæ tuæ, quod nostra peccáta præpédiunt,~* indulgéntia tuæ propitiatiónis accéleret. Qui vivis et regnas.}{english}{\noindent Let us pray. Stir up, we beseech thee, O Lord, thy power, and come; make haste to our aid with thy great might; that, by the help of thy grace, that which is hindered by our sins may be hastened by thy merciful forgiveness. Who livest and reignest.}
\end{document}